% Ce chapitre repose sur :
%   DOC: s12, s13, s20, s24, s30
%   OBS: REL-PAT.H12, REL-PAT.H14, REL-IPP.O03, REL-IPP.R12
%   HYP: absentéisme-par-spécialité, complétude-par-champ, ancienneté-des-créances,
%   HYP: répartition-des-niveaux-de-tri
%   CHI: (aucun)
%   SECTIONS: 5

\chapter{Recommandations}\label{chap:recommandations}

Ce chapitre propose cinq actions. Chacune tient en un bloc de forme fixe, et ce bloc sépare deux
choses que l'on confond ordinairement.

\textbf{Ce qui établit le constat} est ce qui autorise à dire qu'un problème existe. Sur un jeu de
données produit pour ce projet, cela ne peut pas être une valeur affichée par le tableau de bord :
les relations qui gouvernent ces valeurs y ont été injectées, et les relire à la sortie serait
circulaire. Le constat vient donc d'une source extérieure au jeu — un rapport de contrôle, une
publication statistique, une observation au poste — ou bien il n'existe pas, et le rapport le dit.

\textbf{Ce que l'indicateur démontre} est autre chose, et ce n'est pas rien : il démontre que la
chaîne de données sait produire le nombre, à la maille utile, à partir des tables. Le jour où elle
sera branchée sur les données réelles du service, le nombre sera un constat. Aujourd'hui, il est une
capacité.

\textbf{Aucun effort n'est chiffré et aucun effet n'est quantifié.} Un effort en jours et un gain en
points seraient des nombres sans mesure qui les produise, et ce rapport ne s'en autorise pas. Chaque
bloc dit ce que l'action suppose — quel champ existe déjà, quel processus change — et par quelle
mesure on constaterait son effet. L'effet lui-même, lorsqu'il est énoncé, l'est comme une attente.

\section{Sur les rendez-vous}

\textbf{Le constat.} Le délai d'obtention d'un rendez-vous varie fortement d'une spécialité à
l'autre, et une part des créneaux est perdue par des patients qui ne se présentent pas. Le rapport
de la Cour des comptes consacré au centre hospitalier établit l'un et l'autre pour cet établissement
même : il relève des délais moyens de rendez-vous par spécialité, et l'absence de tout état
comparatif entre rendez-vous demandés et consultations réalisées \cite{s20}. L'ordre de grandeur de
l'absentéisme ambulatoire est documenté par ailleurs — une moyenne mondiale de 23,5~\% et une
fourchette de 23 à 33~\% selon les contextes \cite{s24}.

\textbf{Ce que l'indicateur démontre.} Le tableau de bord calcule le délai d'obtention en médiane et
au quatre-vingt-dixième centile par activité, le taux d'absence et le taux d'annulation par
activité. Il démontre que ces quatre grandeurs se produisent depuis le fait des rendez-vous, à la
maille où une décision se prend. Leurs valeurs, en revanche, sont fixées par des paramètres
injectés\convention{absentéisme-par-spécialité} : le classement des spécialités par absentéisme est
imposé, et aucune conclusion sur la force du lien entre délai et absence ne peut en être tirée.

\textbf{L'action.} Rappeler les patients dont le rendez-vous approche, en commençant par les
activités où le délai est le plus long, et ajuster le nombre de créneaux ouverts au taux d'absence
mesuré sur chaque activité une fois la chaîne branchée sur les données réelles.

\textbf{Ce qu'elle suppose.} Le canal existe déjà : la fiche patient porte quatre lignes de
téléphone, une adresse électronique, et deux cases d'avertissement — l'une par message court,
l'autre par courrier électronique\releve{REL-PAT.H12}. Le relevé note que ces cases existent et ne
sont pas exploitées par le service\releve{REL-PAT.H14}. L'action ne demande donc aucun développement
logiciel : elle demande que le champ soit renseigné à l'admission et qu'un envoi soit déclenché.

\textbf{L'effet attendu et sa mesure.} On attend une baisse du taux d'absence sur les activités
ciblées, et une réduction du délai médian sur ces mêmes activités par récupération des créneaux
perdus. La mesure existe : le taux d'absence et le délai d'obtention par activité, comparés avant et
après, à activité constante. La distinction entre absence et annulation, que le système enregistre
par un couple de champs distinct, permet de vérifier que la baisse vient bien des absences et non
d'un report de qualification.

\section{Sur les urgences}

\textbf{Le constat.} Une part importante des passages aux urgences relève d'une consultation
ordinaire. L'avis du Conseil Économique, Social et Environnemental l'établit à l'échelle nationale :
sur les 6~482~185 consultations et soins d'urgence assurés chaque année par le secteur public,
\textbf{64~\% sont des consultations non urgentes} et 10~\% seulement des urgences vitales
\cite{s12}. Une étude conduite dans un hôpital provincial marocain mesure 30,7~\% de consultations
non appropriées sur 410 patients \cite{s13}. Les deux sources concordent sur l'existence du
phénomène ; elles divergent sur son ampleur, l'une comptant les motifs de recours et l'autre le
caractère approprié de la venue.

\textbf{Ce que l'indicateur démontre.} Le tableau de bord calcule la part des passages classés aux
deux niveaux de tri les moins graves, le délai de prise en charge médicale par niveau, la durée de
passage et l'orientation de sortie. Il démontre que la file d'attente se décrit par niveau de
gravité et que le délai s'y mesure. La pyramide de gravité affichée est en revanche
imposée\convention{répartition-des-niveaux-de-tri}, y compris la part relevant d'une consultation
ordinaire : c'est un paramètre, non une observation.

\textbf{L'action.} Ouvrir une filière courte pour les deux niveaux de tri les moins graves, prise en
charge séparément de la file principale, et réorienter en amont vers les établissements de soins de
santé primaires les venues qui relèvent d'eux.

\textbf{Ce qu'elle suppose.} Un tri formalisé à l'accueil, dont le résultat est enregistré : sans
niveau de tri saisi, ni la filière ni sa mesure n'existent. Elle suppose aussi une capacité d'accueil
en amont, que ce rapport ne mesure pas.

\textbf{L'effet attendu et sa mesure.} On attend une baisse du délai de prise en charge des niveaux
les plus graves, la file courte cessant de les précéder. La mesure est le délai de prise en charge
médicale par niveau de tri, en médiane et au quatre-vingt-dixième centile : c'est sur les niveaux
graves qu'il faut la lire, et non sur la moyenne d'ensemble, qu'une filière courte améliore
mécaniquement.

\section{Sur la facturation et le recouvrement}

\textbf{Le constat.} Il est établi pour cet établissement, et par une source extérieure au jeu de
données. La Cour des comptes relève \textbf{50~876~365 dirhams de créances non recouvrées} au
31 décembre 2016, \textbf{57~\% des dossiers non facturés} en 2015, et \textbf{3~477 consultations
payées sur 160~659} aux urgences en 2016 \cite{s20}. Deux problèmes distincts s'y lisent : des
créances qui vieillissent sans être recouvrées, et des épisodes de soin qui ne donnent lieu à aucune
facture.

\textbf{Ce que l'indicateur démontre.} Le tableau de bord calcule le taux de recouvrement des
créances, leur répartition par tranche d'ancienneté, l'aboutissement des relances, le taux
d'encaissement, et le nombre d'épisodes sans facture rattachée par famille et par mois. Il démontre
que la chaîne sait rapprocher un épisode de sa facture et une facture de son encaissement — ce qui
est précisément ce dont l'absence était relevée. Les montants et l'ancienneté affichés sont
imposés\convention{ancienneté-des-créances}, et la priorisation qu'on en tirerait serait circulaire.

\textbf{L'action.} Prioriser les relances par ancienneté décroissante et par type de débiteur, et
produire chaque mois la liste des épisodes sans facture rattachée pour reprise par le service.

\textbf{Ce qu'elle suppose.} Que l'épisode et la facture portent une clé commune exploitable, ce que
le modèle de données démontre possible. Elle suppose surtout un changement de processus et non
d'outil : la créance naît au service d'accueil, où la facture est établie et le règlement encaissé à
la sortie, tandis que le recouvrement relève du pôle administratif — la liste doit donc traverser
cette frontière pour servir.

\textbf{L'effet attendu et sa mesure.} On attend une hausse du taux de recouvrement et une baisse de
la part des créances les plus anciennes. La mesure du rendement de la relance existe — nombre de
relances émises, nombre ayant abouti, et leur rapport — et elle doit être lue avant d'augmenter le
volume des relances plutôt qu'après.

\section{Sur la qualité du fichier patient}

\textbf{Le constat, et il commence par ce qu'on ne sait pas.} Le système n'expose pas quels champs
de la fiche patient sont obligatoires à la saisie et quels champs sont facultatifs. L'observation au
poste ne l'a pas relevé, et le relevé de champs le déclare explicitement : le caractère obligatoire
ou facultatif n'a été relevé sur aucun écran, et n'est donc affirmé nulle part. Recommander de rendre
certains champs obligatoires supposerait de savoir lesquels ne le sont pas. \textbf{C'est le premier
constat, et il fonde la première action.}

Le second constat est établi, lui, et par le système lui-même : une même personne y possède
plusieurs dossiers. L'écran de recherche d'identifiant porte un onglet de recherche pondérée sur
l'index maître des patients\releve{REL-IPP.O03}, et ses résultats s'accompagnent d'une colonne de
probabilité\releve{REL-IPP.R12}. Un système qui offre de chercher un patient par similarité admet
que la recherche exacte ne suffit pas. Le rapport de la Cour des comptes le confirme de l'extérieur,
relevant l'attribution de plusieurs identifiants au même patient et recommandant l'enregistrement de
la pièce d'identité \cite{s20}.

\textbf{Ce que l'indicateur démontre.} Le tableau de bord calcule le taux de renseignement de chaque
colonne de chaque table, le taux de mise en quarantaine, le nombre de groupes de patients partageant
exactement les mêmes valeurs d'identité, et le nombre de grappes formées par le rapprochement
probabiliste. Il démontre que la complétude se mesure colonne par colonne et que l'ampleur du
doublon se chiffre. Le classement des champs par complétude est en revanche
imposé\convention{complétude-par-champ}, et — contrairement aux trois recommandations précédentes —
\textbf{aucune source extérieure ne vient l'établir}. C'est la recommandation la moins appuyée du
chapitre, et il vaut mieux l'écrire que le taire.

\textbf{L'action, en deux temps.} D'abord établir quels champs sont facultatifs, par relevé du
paramétrage du système ou à défaut par mesure de la complétude réelle sur un extrait ; puis rendre
obligatoires ceux qui portent l'identité et le contact. Ensuite, déplacer le contrôle de doublon :
le rapprochement doit s'exécuter au moment où l'identifiant est créé, et non après coup sur un
fichier déjà pollué.

\textbf{Ce qu'elle suppose.} Pour le second temps, presque rien : la fonction existe déjà dans le
produit installé, sous la forme de la recherche pondérée relevée à l'écran. Ce qui manque n'est pas
l'outil mais son insertion dans le geste d'accueil — la recherche doit être obligatoire avant
création, non facultative après.

\textbf{L'effet attendu et sa mesure.} On attend une baisse du nombre de groupes d'identité formés
par collision exacte, et une hausse du taux de renseignement des champs rendus obligatoires. Les deux
mesures existent, et la seconde vérifie la première : un champ rendu obligatoire mais renseigné par
une valeur de remplissage se verrait à la complétude sans améliorer le rapprochement.

\section{Sur les indicateurs que le système ne permet pas de produire}

\textbf{Cette recommandation ne dépend d'aucune donnée générée}, et c'est la seule du chapitre dans
ce cas. Elle porte sur des champs absents, fait qui ne se lit pas dans des valeurs mais dans une
confrontation entre ce que la statistique nationale attend d'un hôpital nommé et ce que la chaîne
peut calculer.

\textbf{Le constat.} Le recueil statistique national publie, pour chaque établissement nommé, une
liste précise d'indicateurs \cite{s30}. Trois d'entre eux ne sont pas calculables, et pour chacun le
champ manquant se nomme.

\begin{center}
\begin{tabular}{@{}p{4.4cm}p{4.6cm}p{5.0cm}@{}}
\toprule
Indicateur attendu & Champ manquant & Ce qu'il débloquerait \\
\midrule
Capacité litière existante & une table de structure de l'établissement, portant les lits installés & la distinction entre lits installés et lits en service, et le suivi de son évolution \\
\addlinespace
Consultations par médecin & un référentiel des médecins affectés & le rapport de l'activité à l'effectif, seul indicateur de productivité que la nomenclature attende \\
\addlinespace
Accouchements & un marqueur d'accouchement sur le séjour & le décompte des accouchements, et le taux de césarienne s'il devenait pertinent \\
\bottomrule
\end{tabular}
\end{center}

\begin{center}
\footnotesize Les trois indicateurs attendus par la statistique nationale que la chaîne ne peut pas
produire, et le champ dont l'absence l'en empêche.
\end{center}

\medskip

\textbf{Une distinction gouverne ce tableau et doit être respectée.} Un indicateur non calculable
faute de champ est un défaut du système d'information, et il appelle une recommandation. Un
indicateur sans objet pour ce site est une propriété de l'établissement, et il n'appelle aucun
reproche au logiciel : les interventions chirurgicales et le taux de césarienne relèvent de cette
seconde catégorie, l'établissement n'exerçant pas d'activité chirurgicale. Confondre les deux
reprocherait au produit une absence de champ là où c'est l'hôpital qui n'a pas l'activité.

\textbf{L'action.} Ajouter les trois champs, en commençant par le référentiel des médecins, dont
dépend le seul indicateur de productivité attendu.

\textbf{Ce qu'elle suppose.} Une intervention sur le paramétrage du produit ou une demande à
l'éditeur, selon que les champs existent et ne sont pas alimentés ou qu'ils n'existent pas — ce que
l'observation conduite ici ne permet pas de trancher.

\textbf{L'effet attendu et sa mesure.} On attend que les trois indicateurs deviennent calculables,
et l'effet se constate sans ambiguïté : la valeur existe ou elle n'existe pas. La confrontation entre
la nomenclature attendue et ce que la chaîne produit, déjà conduite une fois, se refait à l'identique
et donne la réponse.
